\documentclass[10pt]{article}
\usepackage[a4paper,margin=0.84in]{geometry}
\usepackage{amsmath,amssymb,amsthm,mathtools,booktabs,enumitem,xcolor,hyperref}
\hypersetup{colorlinks=true,linkcolor=blue!55!black,citecolor=blue!55!black,urlcolor=blue!55!black}
\setlength{\emergencystretch}{3em}
\setlength{\parindent}{0pt}
\setlength{\parskip}{0.52em}
\setlist{leftmargin=1.65em,itemsep=0.22em,topsep=0.25em}
\sloppy

\newtheorem{theorem}{Theorem}[section]
\newtheorem{proposition}[theorem]{Proposition}
\newtheorem{corollary}[theorem]{Corollary}
\newtheorem{lemma}[theorem]{Lemma}
\theoremstyle{definition}
\newtheorem{definition}[theorem]{Definition}
\newtheorem{convention}[theorem]{Convention}
\newtheorem{remark}[theorem]{Remark}

\newcommand{\Fq}{\mathbf F_q}
\newcommand{\Fbar}{\overline{\mathbf F}_q}
\newcommand{\Ql}{\mathbf Q_\ell}
\newcommand{\Qlb}{\overline{\mathbf Q}_\ell}
\newcommand{\Zl}{\mathbf Z_\ell}
\newcommand{\Qbar}{\overline{\mathbf Q}}
\newcommand{\Tr}{\operatorname{Tr}}
\newcommand{\End}{\operatorname{End}}
\newcommand{\Gal}{\operatorname{Gal}}
\newcommand{\Fix}{\operatorname{Fix}}
\newcommand{\ord}{\operatorname{ord}}
\newcommand{\Spec}{\operatorname{Spec}}
\newcommand{\deps}[1]{\par\smallskip\noindent\textit{Prerequisites.} #1}

\begin{document}
\begin{center}
{\large\bfseries Zeta Functions, Frobenius Determinants, and Weight Separation\par}
\medskip
{\normalsize Determinants, rationality, integrality, and functional equations\par}
\end{center}

\begin{abstract}
Let \(X_0\) be a finite-type scheme over a finite field.  This note develops
the algebraic passage from finite Frobenius orbits and the trace formula to
the cohomological factorization of \(Z(X_0,T)\).  It then separates the
individual factors by weight, descends them to integral polynomials
independent of the auxiliary prime, and derives the raw and normalized
functional equations from perfect Frobenius pairings.  Every intermediate
notion used in the statements is defined explicitly.
\end{abstract}

\tableofcontents

\section{Formal power series at the origin}

Throughout the first four sections, \(R\) is a commutative ring.  Write

\[
 R[[T]]=\varprojlim_{N\geq 0}R[T]/(T^{N+1}).
\]

Thus every infinite identity is coefficientwise, or equivalently an
identity in every finite quotient \(R[T]/(T^{N+1})\).

\begin{definition}[Order and finite-degree agreement]
For \(A=\sum_{n\geq0}a_nT^n\in R[[T]]\), put
\[
 \ord_T(A)=\inf\{n\in\mathbf N:a_n\neq0\}\in
 \mathbf N\cup\{+\infty\}.
\]
For \(A,B\in R[[T]]\) and \(N\geq0\), write
\[
 A\equiv_N B
 \quad\Longleftrightarrow\quad
 A-B\in T^{N+1}R[[T]].
\]
Equivalently, \(A\) and \(B\) have equal coefficients through degree \(N\).
\end{definition}

\begin{definition}[Formal derivative]
The formal derivative is the \(R\)-linear map
\[
 \frac{d}{dT}:R[[T]]\longrightarrow R[[T]],\qquad
 \frac{d}{dT}\left(\sum_{n\geq0}a_nT^n\right)
 =\sum_{n\geq0}(n+1)a_{n+1}T^n.
\]
\end{definition}

\begin{lemma}[Derivative determination]
If \(R\) is a \(\mathbf Q\)-algebra and \(A,B\in R[[T]]\) have the same
constant coefficient and the same formal derivative, then \(A=B\).
\deps{The coefficient formula for the formal derivative and invertibility
of every positive integer in \(R\).}
\end{lemma}

Assume from now until the end of this section that \(R\) is a commutative
\(\mathbf Q\)-algebra.

\begin{definition}[Formal exponential and logarithm]
For \(U\in TR[[T]]\) and \(V\in TR[[T]]\), define
\[
 \operatorname{Exp}(U)=\sum_{r\geq0}\frac{U^r}{r!}
 \in1+TR[[T]],
\]
\[
 \operatorname{Log}(1+V)
 =\sum_{r\geq1}(-1)^{r+1}\frac{V^r}{r}
 \in TR[[T]].
\]
For a fixed coefficient only finitely many powers contribute because
\(\ord_T(U),\ord_T(V)\geq1\).
\end{definition}

\begin{lemma}[Finite coefficient formulas]
Write \(U=\sum_{n\geq1}u_nT^n\) and
\(V=\sum_{n\geq1}v_nT^n\).  For \(N\geq1\),
\[
 [T^N]\operatorname{Exp}(U)
 =\sum_{r=1}^{N}\frac1{r!}
   \sum_{\substack{n_1+\cdots+n_r=N\\n_j\geq1}}
   u_{n_1}\cdots u_{n_r},
\]
\[
 [T^N]\operatorname{Log}(1+V)
 =\sum_{r=1}^{N}\frac{(-1)^{r+1}}r
   \sum_{\substack{n_1+\cdots+n_r=N\\n_j\geq1}}
   v_{n_1}\cdots v_{n_r}.
\]
Both constant coefficients are respectively \(1\) and \(0\).
\deps{The positive \(T\)-adic order of \(U\) and \(V\), and the Cauchy
product formula in \(R[[T]]\).}
\end{lemma}

\begin{lemma}[Formal differential identities]
For \(U\in TR[[T]]\) and \(A\in1+TR[[T]]\),
\[
 \frac{d}{dT}\operatorname{Exp}(U)
 =U'\operatorname{Exp}(U),
 \qquad
 \frac{d}{dT}\operatorname{Log}(A)=A'A^{-1}.
\]
Here \(A^{-1}\) denotes the unique multiplicative inverse of \(A\) in
\(R[[T]]\).
\deps{The finite coefficient formulas, the product rule, and the geometric
inverse of a series with constant coefficient \(1\).}
\end{lemma}

\begin{proposition}[Formal exponential and logarithm]
The maps
\[
 \operatorname{Exp}:(TR[[T]],+)\longrightarrow(1+TR[[T]],\cdot),
 \qquad
 \operatorname{Log}:(1+TR[[T]],\cdot)\longrightarrow(TR[[T]],+)
\]
are inverse group isomorphisms.  In particular,
\[
 \operatorname{Exp}(U+V)=\operatorname{Exp}(U)\operatorname{Exp}(V),
 \qquad
 \operatorname{Log}(AB)=\operatorname{Log}(A)+\operatorname{Log}(B).
\]
\deps{The formal differential identities, the constant coefficients, and
derivative determination.}
\end{proposition}

\begin{lemma}[One inverse Euler factor]
For every integer \(m\geq1\),
\[
 (1-T^m)^{-1}=\sum_{r\geq0}T^{mr},
 \qquad
 \operatorname{Log}(1-T^m)^{-1}
 =\sum_{r\geq1}\frac{T^{mr}}r.
\]
\deps{The geometric-series identity and the formal logarithm of a
multiplicative inverse.}
\end{lemma}

\section{Locally finite products in the T-adic topology}

\begin{definition}[\(T\)-adically locally finite family]
Let \(J\) be a set and let
\(\Phi=(\Phi_j)_{j\in J}\) be a family in \(1+TR[[T]]\).  It is
\emph{\(T\)-adically locally finite} if, for every \(N\geq0\), the active
set
\[
 J_N(\Phi)=\{j\in J:\Phi_j\not\equiv_N1\}
\]
is finite.
\end{definition}

\begin{definition}[Coefficientwise locally finite product]
For a \(T\)-adically locally finite family, define
\[
 \prod_{j\in J}\Phi_j\in1+TR[[T]]
\]
by requiring, for every \(N\),
\[
 \prod_{j\in J}\Phi_j
 \equiv_N\prod_{j\in J_N(\Phi)}\Phi_j.
\]
The product on the right is finite and is computed modulo \(T^{N+1}\).
\end{definition}

\begin{lemma}[Compatibility of finite truncations]
The preceding congruences determine a unique power series.  If
\(N\leq M\), reduction of the level-\(M\) product modulo \(T^{N+1}\)
equals the level-\(N\) product.
\deps{Monotonicity \(J_N(\Phi)\subseteq J_M(\Phi)\) and the fact that every
new factor is congruent to \(1\) modulo \(T^{N+1}\).}
\end{lemma}

\begin{lemma}[Subring preservation]
Let \(R_0\subseteq R\) be a subring.  If every \(\Phi_j\) has coefficients
in \(R_0\), then the locally finite product belongs to
\(1+TR_0[[T]]\).
\deps{The finite truncation definition and coefficientwise injectivity of
\(R_0[[T]]\to R[[T]]\).}
\end{lemma}

\begin{proposition}[Logarithm of a locally finite product]
Assume that \(R\) is a \(\mathbf Q\)-algebra.  For every \(T\)-adically
locally finite family \(\Phi\), the family
\((\operatorname{Log}\Phi_j)_{j\in J}\) is coefficientwise locally finite
and
\[
 \operatorname{Log}\left(\prod_{j\in J}\Phi_j\right)
 =\sum_{j\in J}\operatorname{Log}(\Phi_j).
\]
In degree \(N\), the sum on the right may be restricted to \(J_N(\Phi)\).
\deps{Compatibility of truncations, the finite-product logarithm identity,
and preservation of finite-degree agreement by the formal logarithm.}
\end{proposition}

\section{Periodic permutations and orbit counting}

\begin{definition}[Locally finite periodic permutation]
A locally finite periodic permutation is a pair \((S,F)\) consisting of a
set \(S\) and a bijection \(F:S\to S\) such that:
\begin{enumerate}[label=\textup{(\roman*)}]
\item every \(F\)-orbit in \(S\) is finite;
\item \(\Fix(F^n)=\{s\in S:F^n(s)=s\}\) is finite for every \(n\geq1\).
\end{enumerate}
For an orbit \(O\in S/\langle F\rangle\), its degree is
\(\deg(O)=\#O\).  Put
\[
 a_m=\#\{O\in S/\langle F\rangle:\deg(O)=m\},
 \qquad N_n=\#\Fix(F^n).
\]
\end{definition}

\begin{lemma}[Finiteness of one-degree orbit sets]
For every \(m\geq1\), the set of \(F\)-orbits of degree \(m\) is finite.
\deps{Every degree-\(m\) orbit is contained in the finite set
\(\Fix(F^m)\), and distinct orbits are disjoint.}
\end{lemma}

\begin{lemma}[Fixed points on one orbit]
If \(O\) is an \(F\)-orbit of degree \(m\), then, for every \(n\geq1\),
\[
 \#(O\cap\Fix(F^n))=
 \begin{cases}
  m,&m\mid n,\\
  0,&m\nmid n.
 \end{cases}
\]
\deps{The minimal-period criterion for the action of a cyclic permutation
on one finite orbit.}
\end{lemma}

\begin{proposition}[Fixed points counted orbit by orbit]
For every \(n\geq1\),
\[
 N_n=\sum_{m\mid n}m a_m.
\]
\deps{The fixed-point formula on one orbit, finiteness of \(\Fix(F^n)\),
and finiteness of the divisors of \(n\).}
\end{proposition}

\begin{corollary}[Mobius recovery of orbit counts]
Let \(\mu_{\mathrm{Mob}}:\mathbf N_{>0}\to\{-1,0,1\}\) be the
number-theoretic Mobius function.  Then
\[
 m a_m=\sum_{d\mid m}\mu_{\mathrm{Mob}}(d)N_{m/d},
 \qquad
 a_m=\frac1m\sum_{d\mid m}\mu_{\mathrm{Mob}}(d)N_{m/d}.
\]
In particular, the right-hand numerator is a nonnegative integer divisible
by \(m\).
\deps{The orbit-counting identity and Mobius inversion on the divisibility
poset.}
\end{corollary}

\begin{definition}[Orbit Euler product]
For a locally finite periodic permutation, define
\[
 Z_{\mathrm{orb}}(S,F;T)
 =\prod_{O\in S/\langle F\rangle}
  (1-T^{\deg(O)})^{-1}.
\]
This product is \(T\)-adically locally finite because only finitely many
orbits have degree at most any fixed \(N\).
\end{definition}

\begin{definition}[Counting exponential]
Define
\[
 Z_{\mathrm{cnt}}(S,F;T)
 =\operatorname{Exp}\left(\sum_{n\geq1}N_n\frac{T^n}{n}\right)
 \in1+T\mathbf Q[[T]].
\]
\end{definition}

\begin{theorem}[Orbit Euler product identity]
For every locally finite periodic permutation,
\[
 Z_{\mathrm{orb}}(S,F;T)=Z_{\mathrm{cnt}}(S,F;T).
\]
More precisely, the coefficient of \(T^n\) in the logarithm of the orbit
product is
\[
 \sum_{m\mid n}\frac{a_m}{n/m}=\frac{N_n}{n}.
\]
\deps{The logarithm of a locally finite product, the logarithm of one Euler
factor, and the orbit-counting identity.}
\end{theorem}

\begin{corollary}[Integrality of the counting exponential]
One has
\[
 Z_{\mathrm{cnt}}(S,F;T)\in1+T\mathbf Z[[T]].
\]
\deps{The orbit Euler product identity and subring preservation for locally
finite products of integral Euler factors.}
\end{corollary}

\section{Closed points and the Hasse--Weil series}

Let \(X_0\) be a scheme of finite type over \(\Fq\), let
\(X=X_0\times_{\Fq}\Fbar\), and let
\(F_X:X(\Fbar)\to X(\Fbar)\) denote geometric Frobenius.

\begin{definition}[Point counts]
For \(n\geq1\), define
\[
 X_0(\mathbf F_{q^n})
 =\operatorname{Hom}_{\Fq}(\Spec\mathbf F_{q^n},X_0),
 \qquad
 N_n(X_0)=\#X_0(\mathbf F_{q^n}).
\]
\end{definition}

\begin{definition}[Closed-point degree]
For a closed point \(x\in|X_0|\), its residue field is a finite extension
of \(\Fq\), and
\[
 \deg(x)=[\kappa(x):\Fq].
\]
\end{definition}

\begin{theorem}[Closed-point--orbit dictionary]
Geometric points with the permutation \(F_X\) form a locally finite
periodic permutation.  There are canonical bijections
\[
 X_0(\mathbf F_{q^n})\xrightarrow{\sim}\Fix(F_X^n),
 \qquad
 |X_0|\xrightarrow{\sim}X(\Fbar)/\langle F_X\rangle,
\]
and the second bijection identifies residue degree with orbit cardinality.
Consequently, only finitely many closed points have degree at most a fixed
integer.
\deps{Finite-field descent of geometric points, the description of a closed
point by its geometric lifts, and finiteness of rational points on a
finite-type scheme over a finite field.}
\end{theorem}

\begin{definition}[Hasse--Weil zeta series]
The Hasse--Weil zeta series is
\[
 Z(X_0,T)=
 \operatorname{Exp}\left(\sum_{n\geq1}N_n(X_0)\frac{T^n}{n}\right)
 \in1+T\mathbf Q[[T]].
\]
Its closed-point Euler product is the locally finite product
\[
 \prod_{x\in|X_0|}(1-T^{\deg(x)})^{-1}.
\]
\end{definition}

\begin{corollary}[The two point-counting presentations]
For every finite-type \(\Fq\)-scheme,
\[
 \boxed{
 Z(X_0,T)=
 \prod_{x\in|X_0|}(1-T^{\deg(x)})^{-1}}
 \quad\text{in }1+T\mathbf Z[[T]].
\]
\deps{The closed-point--orbit dictionary, the orbit Euler product identity,
and integrality of the counting exponential.}
\end{corollary}

\section{Finite-dimensional determinants and perfect complexes}

Let \(E\) be a field of characteristic zero.

\begin{definition}[Reciprocal characteristic polynomial]
For a finite-dimensional \(E\)-vector space \(V\) and
\(u\in\End_E(V)\), define
\[
 P_{V,u}(T)=\det(1-Tu\mid V)\in E[T].
\]
After scalar extension to an algebraic closure \(\Omega/E\), write
\[
 P_{V,u}(T)=\prod_{j=1}^{\dim_E V}(1-\alpha_jT).
\]
The multiset \(\{\alpha_j\}\), with algebraic multiplicities, is the
multiset of \emph{reciprocal roots}.  It is the eigenvalue multiset of \(u\).
If \(u\) is invertible, every \(\alpha_j\neq0\),
\(\deg P_{V,u}=\dim_E V\), and
\[
 [T^{\dim V}]P_{V,u}=(-1)^{\dim V}\det(u).
\]
\end{definition}

\begin{lemma}[Finite-dimensional trace--determinant identity]
Let \(V\) be a finite-dimensional \(E\)-vector space and let
\(u\in\End_E(V)\).  Then
\[
 P_{V,u}(T)^{-1}
 =\operatorname{Exp}\left(
   \sum_{n\geq1}\Tr(u^n\mid V)\frac{T^n}{n}\right)
 \quad\text{in }E[[T]].
\]
Equivalently,
\[
 T\frac{d}{dT}\operatorname{Log}P_{V,u}(T)^{-1}
 =\sum_{n\geq1}\Tr(u^n\mid V)T^n.
\]
No semisimplicity assumption is imposed on \(u\).
\deps{The matrix identity
\(\frac{d}{dT}\det(1-Tu)=-\det(1-Tu)
\Tr((1-Tu)^{-1}u)\), the formal resolvent
\((1-Tu)^{-1}=\sum_{n\geq0}T^nu^n\), and derivative determination.}
\end{lemma}

\begin{definition}[Perfect complex over a field]
Let \(D(E)\) be the ordinary derived category of cochain complexes of
\(E\)-vector spaces.  An object \(C\in D(E)\) is \emph{perfect} if it is
isomorphic to a bounded complex of finite-dimensional vector spaces.
Equivalently, \(H^i(C)\) is finite-dimensional for every \(i\) and vanishes
for all but finitely many \(i\).
\end{definition}

\begin{definition}[Supertrace and cohomological zeta factor]
For a perfect complex \(C\) and an endomorphism \(u:C\to C\) in \(D(E)\),
write \(u_i=H^i(u)\) and define
\[
 \operatorname{str}(u\mid C)
 =\sum_i(-1)^i\Tr(u_i\mid H^i(C)),
\]
\[
 Z(C,u;T)
 =\prod_i\det(1-Tu_i\mid H^i(C))^{(-1)^{i+1}}
 \in E(T).
\]
A negative exponent denotes the inverse in \(E(T)\).  Every factor has
constant coefficient \(1\), so \(Z(C,u;T)\) is regular at \(T=0\).
\end{definition}

\begin{proposition}[Additivity on distinguished triangles]
Let \(C'\to C\to C''\to C'[1]\) be a distinguished triangle of perfect
complexes equipped with a commuting endomorphism of the triangle.  Then
for every \(n\geq1\),
\[
 \operatorname{str}(u^n\mid C)
 =\operatorname{str}(u'^n\mid C')+
  \operatorname{str}(u''^n\mid C''),
\]
and
\[
 Z(C,u;T)=Z(C',u';T)Z(C'',u'';T).
\]
Moreover \(Z(C[1],u[1];T)=Z(C,u;T)^{-1}\).
\deps{The long exact cohomology sequence, additivity of trace, and
multiplicativity of characteristic polynomials in short exact sequences.}
\end{proposition}

\begin{theorem}[Perfect-complex determinant--trace identity]
For every perfect \(C\) and every endomorphism \(u\),
\[
 \boxed{
 Z(C,u;T)=\operatorname{Exp}\left(
   \sum_{n\geq1}\operatorname{str}(u^n\mid C)\frac{T^n}{n}
 \right).}
\]
\deps{The finite-dimensional trace--determinant identity in each
cohomological degree and the finite-product law for the formal exponential.}
\end{theorem}

\section{The trace datum and cohomological factorization}

\begin{definition}[Finite cohomological trace datum]
A finite cohomological trace datum over \(E\) consists of:
\begin{enumerate}[label=\textup{(\roman*)}]
\item integers \(N_n\in\mathbf Z\) for every \(n\geq1\);
\item a perfect complex \(C\in D(E)\);
\item an automorphism \(F:C\xrightarrow{\sim}C\);
\item the equalities
\[
 N_n=\operatorname{str}(F^n\mid C)
 \quad\text{in }E
 \qquad(n\geq1).
\]
\end{enumerate}
Put \(V_i=H^i(C)\), \(F_i=H^i(F)\), and
\(P_i(T)=\det(1-TF_i\mid V_i)\).
\end{definition}

\begin{lemma}[Counting logarithm equals trace logarithm]
For a finite cohomological trace datum,
\[
 \sum_{n\geq1}N_n\frac{T^n}{n}
 =\sum_{n\geq1}\left(
   \sum_i(-1)^i\Tr(F_i^n\mid V_i)
  \right)\frac{T^n}{n}
 \quad\text{in }TE[[T]].
\]
\deps{The every-iterate trace equalities and coefficientwise scalar
extension from \(\mathbf Q\) to \(E\).}
\end{lemma}

\begin{proposition}[Cohomological determinant factorization]
For a finite cohomological trace datum,
\[
 \operatorname{Exp}\left(\sum_{n\geq1}N_n\frac{T^n}{n}\right)
 =\prod_iP_i(T)^{(-1)^{i+1}}
 \quad\text{in }E[[T]].
\]
The right side is the Taylor expansion of a rational function in \(E(T)\).
\deps{The equality of logarithms and the perfect-complex
determinant--trace identity.}
\end{proposition}

\begin{corollary}[Cohomological factorization for a variety]
Let \(X_0/\Fq\) be separated and of finite type, let
\[
 C_\ell=R\Gamma_c(X_0\times_{\Fq}\Fbar,\Ql)
 \in D^b_{\mathrm{fd}}(\Ql),
\]
and let \(F\) be geometric Frobenius.  Suppose the
Grothendieck--Lefschetz identities
\[
 N_n(X_0)=\sum_i(-1)^i
 \Tr(F^n\mid H_c^i(X,\Ql))
 \qquad(n\geq1)
\]
hold.  Then, with
\[
 P_{i,\ell}(T)=
 \det(1-TF\mid H_c^i(X,\Ql)),
\]
one has
\[
 Z(X_0,T)=\prod_iP_{i,\ell}(T)^{(-1)^{i+1}}
 \quad\text{in }\Ql[[T]].
\]
In particular the scalar extension of \(Z(X_0,T)\) is rational over
\(\Ql\).
\deps{Perfectness of compactly supported cohomology, Frobenius
invertibility, the trace formula for all iterates, and cohomological
determinant factorization.}
\end{corollary}

\section{Rational germs and Hankel descent}

\begin{definition}[Normalized rational germ]
Let \(K\) be a field.  A normalized polynomial fraction is a pair
\((A,B)\in K[T]^2\) with \(B(0)=1\).  Two pairs are equivalent when
\[
 (A,B)\sim(A',B')
 \quad\Longleftrightarrow\quad AB'=A'B.
\]
The resulting equivalence class is a \emph{normalized rational germ at
zero}.  Its rational-function value is \(A/B\in K(T)\), and its Taylor
series is
\[
 \tau_K[A,B]=A(T)B(T)^{-1}\in K[[T]].
\]
Write \(K(T)_0\) for the subring of rational functions regular at \(T=0\).
\end{definition}

\begin{lemma}[Faithfulness of the Taylor map]
If two normalized rational germs have the same Taylor series, then they are
equal.  Equivalently, Taylor expansion gives an injective map
\[
 K(T)_0\hookrightarrow K[[T]].
\]
\deps{Invertibility of denominator series with constant coefficient \(1\),
cross multiplication, and injectivity of \(K[T]\to K[[T]]\).}
\end{lemma}

\begin{definition}[Shift span and Hankel rank]
For a sequence \(a=(a_n)_{n\geq0}\in K^{\mathbf N}\), define the left
shift \(Sa\) by \((Sa)_n=a_{n+1}\), and put
\[
 \operatorname{Sh}_K(a)
 =\operatorname{span}_K\{S^ma:m\geq0\}\subseteq K^{\mathbf N}.
\]
The Hankel matrix is \(H(a)=(a_{r+s})_{r,s\geq0}\).  It has rank at most
\(h\) if every \((h+1)\times(h+1)\) minor
\[
 \det(a_{r_x+s_y})_{0\leq x,y\leq h}
\]
vanishes for all choices of nonnegative indices \(r_x,s_y\).
\end{definition}

\begin{definition}[Eventual recurrence]
An eventual recurrence of order \(d\) for \(a\) is a tuple
\((c_1,\ldots,c_d;N)\) in \(K^d\times\mathbf N\) satisfying
\[
 a_n+c_1a_{n-1}+\cdots+c_da_{n-d}=0
 \qquad(n\geq\max\{N,d\}).
\]
\end{definition}

\begin{lemma}[A rational germ supplies a recurrence]
Let \(f(T)=\sum_{n\geq0}a_nT^n\) be the Taylor series of \(A(T)/B(T)\),
where
\[
 B(T)=1+c_1T+\cdots+c_dT^d.
\]
Then
\[
 a_n+c_1a_{n-1}+\cdots+c_da_{n-d}=0
 \]
whenever \(n>\deg A\) and \(n\geq d\).
\deps{The coefficient identity \(B(T)f(T)=A(T)\).}
\end{lemma}

\begin{lemma}[A recurrence gives finite shift span]
If \(a\) satisfies an eventual recurrence, then
\(\operatorname{Sh}_K(a)\) is finite-dimensional.
\deps{The recurrence expresses every sufficiently late shift in terms of
finitely many earlier shifts and finitely many coordinate vectors supported
before the recurrence threshold.}
\end{lemma}

\begin{lemma}[Finite coordinates detect independent sequences]
Let \(v_1,\ldots,v_r\in K^{\mathbf N}\) be linearly independent.  There
exist indices \(n_1,\ldots,n_r\in\mathbf N\) such that
\[
 \det(v_j(n_i))_{1\leq i,j\leq r}\neq0.
\]
\deps{Finite Gaussian elimination using the first coordinate at which a
nonzero residual sequence does not vanish.}
\end{lemma}

\begin{proposition}[Hankel rank equals shift-span dimension]
For \(a\in K^{\mathbf N}\) and \(h\geq0\), the following are equivalent:
\begin{enumerate}[label=\textup{(\roman*)}]
\item every \((h+1)\times(h+1)\) Hankel minor vanishes;
\item \(\dim_K\operatorname{Sh}_K(a)\leq h\).
\end{enumerate}
\deps{The columns of a Hankel matrix are the shifts of \(a\), and finite
coordinates detect every finite independent family of shifts.}
\end{proposition}

\begin{proposition}[Finite shift span supplies a rational germ]
Assume \(W=\operatorname{Sh}_K(a)\) is finite-dimensional of dimension
\(r\).  The shift restricts to an endomorphism \(S_W:W\to W\).  Write
\[
 \det(U-S_W)=U^r+c_{r-1}U^{r-1}+\cdots+c_0
\]
and set
\[
 B_a(T)=1+c_{r-1}T+\cdots+c_0T^r.
\]
If \(f_a(T)=\sum_{n\geq0}a_nT^n\), then
\(B_a(T)f_a(T)\) is a polynomial \(A_a(T)\).  If \(r=0\), then
\(a=0\) and \(A_a=0\); if \(r>0\), then \(\deg A_a<r\).  In both cases,
\[
 f_a(T)=\tau_K[A_a,B_a].
\]
\deps{Cayley--Hamilton for the shift endomorphism and coefficientwise
multiplication by \(B_a(T)\).}
\end{proposition}

\begin{lemma}[Finite determinants reflect scalar extension]
Let \(K\hookrightarrow L\) be an injective field homomorphism.  For every
finite square matrix \(M\) over \(K\),
\[
 \det(M)=0
 \quad\Longleftrightarrow\quad
 \det(M\otimes_KL)=0.
\]
Consequently, every finite Hankel-rank bound on the scalar extension of a
sequence over \(K\) descends to the original sequence.
\deps{The determinant is a finite polynomial in the matrix entries, and an
injective field homomorphism reflects zero.}
\end{lemma}

\begin{theorem}[Hankel rationality descent]
For every field extension \(K\subseteq L\),
\[
 \boxed{K[[T]]\cap L(T)_0=K(T)_0}
\]
inside \(L[[T]]\).  Explicitly, if \(f\in K[[T]]\) becomes the Taylor
series of a normalized rational germ over \(L\), then its coefficient
sequence has finite shift span over \(K\), and the shift characteristic
polynomial constructs a normalized rational germ over \(K\) with Taylor
series \(f\).
\deps{A rational germ supplies an eventual recurrence, a recurrence gives
a finite Hankel-rank bound, finite determinants reflect scalar extension,
and finite shift span supplies a rational germ.}
\end{theorem}

\begin{corollary}[Rationality over the counting field]
If \(Z(T)\in\mathbf Q[[T]]\) becomes rational after coefficientwise
extension to a characteristic-zero field \(E\), then
\(Z(T)\in\mathbf Q(T)_0\).
\deps{Hankel rationality descent for \(\mathbf Q\subseteq E\).}
\end{corollary}

\section{Integral Fatou and normalized Gauss descent}

\begin{definition}[Integral shift lattice]
Let \(a\in\mathbf Z^{\mathbf N}\) have finite-dimensional rational shift
span \(W=\operatorname{Sh}_{\mathbf Q}(a)\).  Define
\[
 M(a)=\sum_{m\geq0}\mathbf Z\,S^ma\subseteq W.
\]
This is a shift-stable torsion-free abelian group and
\(M(a)\otimes_{\mathbf Z}\mathbf Q=W\).
\end{definition}

\begin{lemma}[Finite-coordinate lattice]
There is a finite set \(J\subset\mathbf N\) such that coordinate evaluation
\[
 \operatorname{ev}_J:W\longrightarrow\mathbf Q^J
\]
is injective.  Its restriction maps \(M(a)\) injectively into
\(\mathbf Z^J\).  Hence \(M(a)\) is a finite free abelian group.
\deps{Finite coordinates detect a basis of \(W\), every element of
\(M(a)\) is an integral sequence, and every subgroup of a finite free
abelian group is finite free.}
\end{lemma}

\begin{lemma}[Integral Cayley--Hamilton recurrence]
The shift endomorphism \(S:M(a)\to M(a)\) has a monic characteristic
polynomial
\[
 C(U)=U^r+c_{r-1}U^{r-1}+\cdots+c_0\in\mathbf Z[U].
\]
For every \(n\geq0\),
\[
 a_{n+r}+c_{r-1}a_{n+r-1}+\cdots+c_0a_n=0.
\]
\deps{Finite freeness of the integral shift lattice, the integral matrix of
its shift, and Cayley--Hamilton over \(\mathbf Z\).}
\end{lemma}

\begin{proposition}[Integral rational presentation]
Let \(f(T)=\sum_{n\geq0}a_nT^n\in\mathbf Z[[T]]\) be rational over
\(\mathbf Q\).  Then there exist \(A_0,B_0\in\mathbf Z[T]\) with
\(B_0(0)=1\) and \(A_0(0)=a_0\) such that
\[
 B_0(T)f(T)=A_0(T)
 \quad\text{in }\mathbf Z[[T]].
\]
One may take
\[
 B_0(T)=1+c_{r-1}T+\cdots+c_0T^r
\]
from the integral shift characteristic polynomial.
\deps{Hankel rationality, the integral Cayley--Hamilton recurrence, and
coefficientwise multiplication of the recurrence by the series \(f\).}
\end{proposition}

\begin{definition}[Primitive polynomial and content]
For a nonzero polynomial \(P\in\mathbf Z[T]\), its content
\(\operatorname{cont}(P)\) is the positive greatest common divisor of its
coefficients.  It is \emph{primitive} if its content is \(1\).  A nonzero
rational polynomial has a unique nonzero rational multiple that is primitive
with
positive leading coefficient; this is its primitive integral
normalization.
\end{definition}

\begin{lemma}[Gauss product and divisibility]
The product of primitive polynomials in \(\mathbf Z[T]\) is primitive.  If
\(P,Q\in\mathbf Z[T]\), \(P\) is primitive, and \(P\) divides \(Q\) in
\(\mathbf Q[T]\), then \(P\) divides \(Q\) in \(\mathbf Z[T]\).
\deps{Multiplicativity of polynomial content and unique factorization in
\(\mathbf Z[T]\).}
\end{lemma}

\begin{lemma}[Monic rational factors of an integral monic polynomial]
Let \(C\in\mathbf Z[U]\) be monic.  If
\[
 C(U)=R(U)S(U)
\]
for monic polynomials \(R,S\in\mathbf Q[U]\), then
\(R,S\in\mathbf Z[U]\).
\deps{Primitive integral normalization, multiplicativity of content, and
Gauss divisibility in \(\mathbf Z[U]\).}
\end{lemma}

\begin{lemma}[Normalized Gauss factor lemma]
Let \(A\in\mathbf Z[T]\) satisfy \(A(0)=1\).  If
\[
 A(T)=P(T)Q(T),\qquad P,Q\in\mathbf Q[T],
 \qquad P(0)=Q(0)=1,
\]
then \(P,Q\in\mathbf Z[T]\).
\deps{The monic rational-factor lemma for the reversed polynomials
\(T^{\deg A}A(T^{-1})\),
\(T^{\deg P}P(T^{-1})\), and
\(T^{\deg Q}Q(T^{-1})\).}
\end{lemma}

\begin{definition}[Coprime polynomials and normalized gcd]
Let \(K\) be a field.  Polynomials \(P,Q\in K[T]\), not both zero, are
\emph{coprime} if
\[
 (P,Q)=K[T],
\]
or equivalently if every common divisor of \(P\) and \(Q\) is a unit.  If
\(P(0)=Q(0)=1\), their \emph{normalized greatest common divisor} is the
unique polynomial \(G\in K[T]\) such that
\[
 G(0)=1,
 \qquad
 (P,Q)=G K[T].
\]
\end{definition}

\begin{proposition}[Normalized gcd reduction over the integers]
Let \(A_0,B_0\in\mathbf Z[T]\) satisfy
\(A_0(0)=B_0(0)=1\), and let \(G\in\mathbf Q[T]\) be their normalized
greatest common divisor in \(\mathbf Q[T]\).  Define
\(A,B\in\mathbf Q[T]\) by
\[
 A_0=GA,
 \qquad
 B_0=GB.
\]
Then
\[
 G,A,B\in\mathbf Z[T],
 \qquad
 G(0)=A(0)=B(0)=1,
\]
the pair \((A,B)\) is coprime in \(\mathbf Q[T]\), and
\[
 \frac{A_0}{B_0}=\frac{A}{B}
 \quad\text{in }\mathbf Q(T).
\]
\deps{Existence and uniqueness of a normalized generator in the principal
ideal domain \(\mathbf Q[T]\), and the normalized Gauss factor lemma applied
to \(A_0=GA\) and \(B_0=GB\).}
\end{proposition}

\begin{lemma}[Uniqueness of normalized reduced fractions]
Let \(K\) be a field.  Suppose \(A,B,O,E\in K[T]\) have constant
coefficient \(1\), both pairs \((A,B)\) and \((O,E)\) are coprime, and
\(OB=EA\).  Then
\[
 O=A,
 \qquad
 E=B.
\]
\deps{Euclid's lemma gives the two divisibilities in each direction, and
the constant coefficient \(1\) removes the remaining scalar units.}
\end{lemma}

\begin{theorem}[Integral Fatou theorem]
Let
\[
 f(T)\in1+T\mathbf Z[[T]]\cap\mathbf Q(T).
\]
There is a unique pair of polynomials \(A,B\in\mathbf Z[T]\), coprime in
\(\mathbf Q[T]\),
satisfying
\[
 A(0)=B(0)=1,
 \qquad f(T)=\frac{A(T)}{B(T)}.
\]
This is the normalized reduced integral presentation of \(f\).
\deps{The integral rational presentation, equality of its numerator and
denominator constant coefficients, normalized gcd reduction over the
integers, and uniqueness of normalized reduced fractions over
\(\mathbf Q[T]\).}
\end{theorem}

\begin{definition}[Reversal and reciprocal roots]
Let \(P(T)=1+c_1T+\cdots+c_bT^b\in\mathbf Z[T]\) have degree \(b\).  Its
reversal is
\[
 P^{\mathrm{rev}}(U)=U^bP(U^{-1})
 =U^b+c_1U^{b-1}+\cdots+c_b\in\mathbf Z[U].
\]
For a field \(\Omega\) containing \(\mathbf Q\), a scalar
\(\alpha\in\Omega^\times\) is a \emph{reciprocal root} of \(P\) if
\(P(\alpha^{-1})=0\).
\end{definition}

\begin{lemma}[Reciprocal roots and reversal]
For \(P\) as above, the polynomial \(P^{\mathrm{rev}}\) is monic.  For every
field \(\Omega\) containing \(\mathbf Q\) and every
\(\alpha\in\Omega^\times\),
\[
 P(\alpha^{-1})=0
 \quad\Longleftrightarrow\quad
 P^{\mathrm{rev}}(\alpha)=0.
\]
\deps{The defining identity \(P^{\mathrm{rev}}(U)=U^bP(U^{-1})\).}
\end{lemma}

\begin{corollary}[Integral reciprocal roots]
Every reciprocal root of a polynomial
\(P(T)=1+c_1T+\cdots+c_bT^b\in\mathbf Z[T]\) of degree \(b\) is an
algebraic integer.
\deps{The monic integral polynomial \(P^{\mathrm{rev}}\) and the
reciprocal-root equivalence.}
\end{corollary}

\section{Universal weights and weight strata}

Fix an integer \(q>1\).

\begin{definition}[Universal complex-realization field]
A universal complex-realization field is an algebraically closed field
\(\Omega\) of characteristic zero, containing \(\mathbf Q\), such that
\(\operatorname{Isom}_{\mathbf Q}(\Omega,\mathbf C)\) is nonempty and every
embedding of a number field contained in \(\Omega\) into \(\mathbf C\)
extends to an element of this set.
\end{definition}

\begin{proposition}[The algebraic closure of an l-adic field is universal]
For every prime \(\ell\), an abstract algebraic closure \(\Qlb\) of
\(\Ql\), regarded only as a field containing \(\mathbf Q\), is a universal
complex-realization field.  In particular, every embedding
\(K\hookrightarrow\mathbf C\) of a number field
\(K\subseteq\Qlb\) extends to a \(\mathbf Q\)-field isomorphism
\[
 \Qlb\xrightarrow{\sim}\mathbf C.
\]
\deps{Extension of embeddings to algebraic closures, transcendence bases,
the equality of the transcendence degrees of \(\Qlb\) and \(\mathbf C\)
over \(\Qbar\), and extension of an isomorphism between transcendence-basis
fields to their algebraic closures.}
\end{proposition}

\begin{definition}[Universal purity of a scalar]
For a universal complex-realization field \(\Omega\), a nonzero
\(\alpha\in\Omega\) is \emph{universally pure of weight \(w\in\mathbf R\)}
if
\[
 |\iota(\alpha)|=q^{w/2}
 \qquad
 \text{for every }\iota\in
 \operatorname{Isom}_{\mathbf Q}(\Omega,\mathbf C).
\]
This condition does not include algebraicity of \(\alpha\) over
\(\mathbf Q\).
\end{definition}

\begin{definition}[Pure polynomial]
A polynomial \(P(T)\in\Omega[T]\) with \(P(0)=1\) is \emph{pure of weight
\(w\)} if all of its reciprocal roots, counted with multiplicity, are
nonzero and universally pure of weight \(w\).
\end{definition}

\begin{lemma}[Uniqueness of a scalar weight]
If a nonzero scalar is universally pure of weights \(w\) and \(w'\), then
\(w=w'\).
\deps{For one complex realization, the equality
\(q^{w/2}=q^{w'/2}\) and injectivity of the real exponential with base
\(q>1\).}
\end{lemma}

\begin{lemma}[Automorphism invariance of universal weight]
If \(\alpha\in\Omega^\times\) is universally pure of weight \(w\) and
\(\sigma\in\operatorname{Aut}_{\mathbf Q}(\Omega)\), then
\(\sigma(\alpha)\) is universally pure of weight \(w\).
\deps{Precomposition of every complex realization with \(\sigma\) is again
a complex realization.}
\end{lemma}

\begin{lemma}[Classical all-embeddings weight]
Let \(\alpha\in\Omega\) be algebraic over \(\mathbf Q\) and universally
pure of weight \(w\).  For every embedding
\(\tau:\mathbf Q(\alpha)\hookrightarrow\mathbf C\),
\[
 |\tau(\alpha)|=q^{w/2}.
\]
\deps{The extension property in the definition of a universal
complex-realization field.}
\end{lemma}

\begin{definition}[Galois-stable weighted root multiset]
Let \(R(T)\in\mathbf Q[T]\) satisfy \(R(0)=1\), and let
\(\mathcal R(R)\) be its reciprocal-root multiset in \(\Qbar\).  For a
real number \(w\), define
\[
 \mathcal R_w(R)=\{\alpha\in\mathcal R(R):
 |\tau(\alpha)|=q^{w/2}
 \text{ for every }\tau:\mathbf Q(\alpha)\hookrightarrow\mathbf C\}.
\]
Multiplicities are inherited from \(R\).
\end{definition}

\begin{proposition}[Galois stability of a weight stratum]
The multiset \(\mathcal R_w(R)\) is stable, with multiplicities, under
\(\Gal(\Qbar/\mathbf Q)\).  Consequently
\[
 R_w(T)=\prod_{\alpha\in\mathcal R_w(R)}(1-\alpha T)
\]
belongs to \(\mathbf Q[T]\), has constant coefficient \(1\), and is the
unique normalized factor of \(R\) with exactly this reciprocal-root
multiset.
\deps{A rational polynomial has a Galois-stable root multiset, the
all-embeddings modulus condition is invariant under Galois conjugation, and
elementary symmetric functions determine a monic normalized factor.}
\end{proposition}

\section{No cancellation between different weights}

Let \(r\geq0\), let \(\Omega\) be a universal complex-realization field,
and let
\[
 P_i(T)\in\Omega[T],\qquad 0\leq i\leq r,
\]
have constant coefficient \(1\).  Define the parity products
\[
 O(T)=\prod_{\substack{0\leq i\leq r\\i\text{ odd}}}P_i(T),
 \qquad
 E(T)=\prod_{\substack{0\leq i\leq r\\i\text{ even}}}P_i(T).
\]

\begin{lemma}[Different weights have disjoint reciprocal roots]
Assume \(P_i\) is pure of weight \(i\) for every \(i\).  If \(i\neq j\),
then \(P_i\) and \(P_j\) have no common reciprocal root in \(\Omega\).
\deps{A common reciprocal root would be universally pure of both weights,
contradicting uniqueness of scalar weight.}
\end{lemma}

\begin{corollary}[Coprimality of the parity products]
Under the same purity hypothesis, the polynomials \(O\) and \(E\) are
coprime in \(\Omega[T]\).
\deps{Different weights have disjoint roots, and two polynomials over an
algebraically closed field are coprime exactly when they have no common
root.}
\end{corollary}

\begin{lemma}[Cross-product identity]
Let \(A,B\in\mathbf Q[T]\) satisfy \(B(0)=1\), and suppose
\[
 \frac{O(T)}{E(T)}=\frac{A(T)}{B(T)}
 \quad\text{in }\Omega(T).
\]
Then
\[
 O(T)B(T)=E(T)A(T)
 \quad\text{in }\Omega[T].
\]
\deps{Both denominators are nonzero because their constant coefficients are
\(1\), so equality in the fraction field may be cross-multiplied.}
\end{lemma}

\begin{proposition}[Identification of the two parity factors]
Assume every \(P_i\) is pure of weight \(i\).  Let
\(A,B\in\mathbf Q[T]\) be coprime with \(A(0)=B(0)=1\), and suppose
\[
 \frac{O(T)}{E(T)}=\frac{A(T)}{B(T)}.
\]
Then, after coefficient extension to \(\Omega\),
\[
 \boxed{O(T)=A(T),\qquad E(T)=B(T).}
\]
Thus no reciprocal root cancels between odd and even cohomological degrees.
\deps{Coprimality of the parity products, the cross-product identity, and
uniqueness of normalized reduced fractions.}
\end{proposition}

\begin{corollary}[Algebraicity of every reciprocal root]
Under the hypotheses of the preceding proposition, every reciprocal root
of every \(P_i\) is algebraic over \(\mathbf Q\).
\deps{Each \(P_i\) divides the scalar extension of \(A\) or \(B\), whose
roots are algebraic over \(\mathbf Q\).}
\end{corollary}

\section{Recovery of individual factors}

Continue with the hypotheses of the preceding section, and assume now that
\(A,B\in\mathbf Z[T]\).

\begin{definition}[Complete weight factor]
For \(0\leq i\leq r\), put
\[
 R_i(T)=
 \begin{cases}
  A(T),&i\text{ odd},\\
  B(T),&i\text{ even}.
 \end{cases}
\]
The \emph{complete weight-\(i\) factor} of the reduced fraction \(A/B\) is
\[
 P_i^{\mathrm{rec}}(T)
 =\prod_{\alpha\in\mathcal R_i(R_i)}(1-\alpha T)
 \in\mathbf Q[T].
\]
\end{definition}

\begin{lemma}[Root multiset of one cohomological degree]
For every \(i\), the reciprocal-root multiset of \(P_i\) is exactly the
weight-\(i\) stratum \(\mathcal R_i(R_i)\).
\deps{The parity product equals \(R_i\), different degree factors have
disjoint weights, and every root of \(R_i\) occurs in one factor of the
corresponding parity.}
\end{lemma}

\begin{proposition}[Unique recovery over the coefficient field]
For every \(i\),
\[
 P_i(T)=P_i^{\mathrm{rec}}(T)
 \quad\text{in }\Omega[T].
\]
In particular \(P_i\) is the scalar extension of a unique polynomial in
\(\mathbf Q[T]\), determined by \(A/B\), \(q\), and the integer \(i\).
\deps{Galois stability of a weight stratum, equality of root multisets,
and the normalization \(P_i(0)=1\).}
\end{proposition}

\begin{proposition}[Integrality of an individual factor]
The recovered polynomial \(P_i^{\mathrm{rec}}(T)\) belongs to
\(\mathbf Z[T]\).
\deps{It is a normalized rational factor of \(A\) or \(B\), and the
normalized Gauss factor lemma applies to that integral parity polynomial
and its normalized complementary factor.}
\end{proposition}

\begin{corollary}[Arithmetic properties of reciprocal roots]
Every reciprocal root \(\alpha\) of \(P_i^{\mathrm{rec}}\) is an algebraic
integer, its Galois conjugates occur with the same multiplicities, and
\[
 |\tau(\alpha)|=q^{i/2}
\]
for every embedding \(\tau:\mathbf Q(\alpha)\hookrightarrow\mathbf C\).
\deps{Integrality of the individual factor, the integral reciprocal-root
lemma, Galois stability of the weight stratum, and classical
all-embeddings weight.}
\end{corollary}

\section{Compatible realizations and independence of the auxiliary prime}

\begin{definition}[Compatible weighted realization family]
Let \(\Lambda\) be a nonempty index set, and fix a base index
\(\lambda_*\in\Lambda\).  For each \(\lambda\in\Lambda\), let
\(\Omega_\lambda\) be a universal complex-realization field and let
\[
 (V_{i,\lambda},F_{i,\lambda})_{0\leq i\leq r}
\]
be finite-dimensional \(\Omega_\lambda\)-vector spaces with automorphisms.
Put
\[
 P_{i,\lambda}(T)
 =\det(1-TF_{i,\lambda}\mid V_{i,\lambda}).
\]
The family is \emph{compatible through a common integral zeta series} if:
\begin{enumerate}[label=\textup{(\roman*)}]
\item every \(P_{i,\lambda}\) is universally pure of weight \(i\);
\item there is a single
\(Z(T)\in1+T\mathbf Z[[T]]\cap\mathbf Q(T)\) such that, for every
\(\lambda\),
\[
 Z(T)=\prod_{i=0}^{r}
 P_{i,\lambda}(T)^{(-1)^{i+1}}
 \quad\text{in }\Omega_\lambda[[T]].
\]
\end{enumerate}
\end{definition}

\begin{lemma}[The reduced fraction is realization-independent]
Let \(Z=A/B\) be the normalized reduced integral presentation furnished by
Integral Fatou.  The pair \((A,B)\) depends only on the coefficient sequence
of \(Z\), and not on \(\lambda\).
\deps{Existence and uniqueness in the Integral Fatou theorem.}
\end{lemma}

\begin{theorem}[Individual-factor independence]
For every \(i\), there is a unique polynomial
\[
 P_i(T)\in\mathbf Z[T],\qquad P_i(0)=1,
\]
whose scalar extension to every \(\Omega_\lambda\) is
\(P_{i,\lambda}(T)\).  It is the complete weight-\(i\) factor of \(A\) for
odd \(i\), and of \(B\) for even \(i\).
\deps{The realization-independent reduced fraction, no cancellation,
unique recovery by the complete weight stratum, and integrality of each
recovered factor.}
\end{theorem}

\begin{corollary}[Independence of cohomological dimensions]
For every \(i\),
\[
 b_i:=\deg P_i=\dim_{\Omega_\lambda}V_{i,\lambda}
\]
is independent of \(\lambda\).
\deps{An automorphism has nonzero determinant, so
\(\deg\det(1-TF)=\dim V\), and scalar extension preserves polynomial
degree.}
\end{corollary}

\begin{corollary}[Complete arithmetic recovery]
The common point-counting series and the distinct weights determine, without
comparison maps between different realizations:
\begin{enumerate}[label=\textup{(\roman*)}]
\item every integral polynomial \(P_i(T)\);
\item every Betti number \(b_i=\deg P_i\);
\item the Galois-stable reciprocal-root multiset in degree \(i\);
\item algebraic integrality and the modulus \(q^{i/2}\) of every conjugate.
\end{enumerate}
\deps{Individual-factor independence, independence of dimensions, and the
arithmetic properties of reciprocal roots.}
\end{corollary}

\section{Perfect similitude pairings}

Let \(E\) be a characteristic-zero field.

\begin{definition}[Perfect similitude pairing]
Let \(V,W\) be finite-dimensional \(E\)-vector spaces.  A bilinear pairing
\[
 B:V\times W\longrightarrow E
\]
is \emph{perfect} if the adjoint map
\[
 \phi_B:W\longrightarrow V^\vee,
 \qquad \phi_B(w)(v)=B(v,w),
\]
is an isomorphism.  Given automorphisms \(F_V,F_W\) and
\(\mu\in E^\times\), it is a \emph{similitude pairing of multiplier
\(\mu\)} if
\[
 B(F_Vv,F_Ww)=\mu B(v,w)
 \qquad(v\in V,\ w\in W).
\]
\end{definition}

\begin{lemma}[Adjoint conjugacy]
Let \(E\) be a field, let \(V,W\) be finite-dimensional \(E\)-vector
spaces, let \(F_V\in\operatorname{Aut}_E(V)\),
\(F_W\in\operatorname{Aut}_E(W)\), and let \(\mu\in E^\times\).  Let
\(B:V\times W\to E\) be perfect and satisfy
\[
 B(F_Vv,F_Ww)=\mu B(v,w)
 \qquad(v\in V,\ w\in W).
\]
Define \(\phi_B:W\to V^\vee\) by \(\phi_B(w)(v)=B(v,w)\).  Then
\[
 \phi_BF_W
 =\mu(F_V^{-1})^\vee\phi_B,
\]
so \(F_W\) is conjugate to
\(\mu(F_V^{-1})^\vee\).
\deps{Perfectness of \(B\), the similitude identity, and the definition of
the dual endomorphism.}
\end{lemma}

\begin{proposition}[Determinants in a similitude pair]
Let \(E\) be a field, let \(V,W\) be finite-dimensional \(E\)-vector
spaces, let \(F_V\in\operatorname{Aut}_E(V)\),
\(F_W\in\operatorname{Aut}_E(W)\), and let \(\mu\in E^\times\).  Let
\(B:V\times W\to E\) be perfect and satisfy
\[
 B(F_Vv,F_Ww)=\mu B(v,w)
 \qquad(v\in V,\ w\in W).
\]
Put
\[
 b=\dim_EV=\dim_EW,\qquad
 D_V=\det(F_V),\qquad D_W=\det(F_W).
\]
Then
\[
 D_VD_W=\mu^b.
\]
\deps{Adjoint conjugacy, determinant invariance under conjugation,
\(\det(F^\vee)=\det(F)\), and
\(\det(\mu F^{-1})=\mu^b\det(F)^{-1}\).}
\end{proposition}

\begin{proposition}[Reciprocal eigenvalue pairing]
Let \(E\) be a field, let \(V,W\) be finite-dimensional \(E\)-vector
spaces, let \(F_V\in\operatorname{Aut}_E(V)\),
\(F_W\in\operatorname{Aut}_E(W)\), and let \(\mu\in E^\times\).  Let
\(B:V\times W\to E\) be perfect and satisfy
\[
 B(F_Vv,F_Ww)=\mu B(v,w)
 \qquad(v\in V,\ w\in W).
\]
Let \(\Omega\) be an algebraic closure of \(E\), and put
\(b=\dim_EV=\dim_EW\).  If
\(\alpha_1,\ldots,\alpha_b\in\Omega\) are the eigenvalues of
\(F_V\otimes_E\Omega\), counted with algebraic multiplicity, then the
eigenvalues of \(F_W\otimes_E\Omega\), counted with algebraic multiplicity,
are
\[
 \frac{\mu}{\alpha_1},\ldots,\frac{\mu}{\alpha_b}.
\]
No semisimplicity assumption is required.
\deps{Adjoint conjugacy and the characteristic-polynomial identities for
dual, inverse, and scalar-multiple endomorphisms.}
\end{proposition}

\begin{theorem}[Degreewise reciprocal-polynomial identities]
Let \(E\) be a field, let \(V,W\) be finite-dimensional \(E\)-vector
spaces, let \(F_V\in\operatorname{Aut}_E(V)\),
\(F_W\in\operatorname{Aut}_E(W)\), and let \(\mu\in E^\times\).  Let
\(B:V\times W\to E\) be perfect and satisfy
\[
 B(F_Vv,F_Ww)=\mu B(v,w)
 \qquad(v\in V,\ w\in W).
\]
Put
\[
 b=\dim_EV=\dim_EW,\qquad D_V=\det(F_V),
\]
and define
\[
 P_V(T)=\det(1-TF_V\mid V),\qquad
 P_W(T)=\det(1-TF_W\mid W).
\]
Then, in \(E[T,T^{-1}]\),
\[
 P_W(\mu^{-1}T^{-1})
 =(-1)^bT^{-b}D_V^{-1}P_V(T),
\]
or equivalently
\[
 P_W(T)=(-1)^b\mu^bD_V^{-1}T^b
 P_V(\mu^{-1}T^{-1}).
\]
\deps{Reciprocal eigenvalue pairing and the product expression
\(P_V(T)=\prod_j(1-\alpha_jT)\), or equivalently adjoint conjugacy and
determinant functoriality over the Laurent polynomial ring.}
\end{theorem}

\section{The middle pairing and its determinant sign}

\begin{definition}[Symmetry type]
Let \(V\) be a finite-dimensional \(E\)-vector space, let
\(B:V\times V\to E\) be a bilinear pairing, and let
\(s\in\{1,-1\}\).  The pairing \(B\) has \emph{swap sign} \(s\) if
\[
 B(y,x)=sB(x,y)
 \qquad(x,y\in V).
\]
The cases \(s=1\) and \(s=-1\) are respectively symmetric and alternating
because \(E\) has characteristic zero.
\end{definition}

\begin{lemma}[Even dimension of a perfect alternating space]
Let \(V\) be a finite-dimensional \(E\)-vector space and let
\(B:V\times V\to E\) be a perfect alternating pairing.  Then
\(\dim_EV\) is even.
\deps{The determinant of an odd-size skew-symmetric matrix is zero, whereas
perfectness makes the pairing matrix invertible.}
\end{lemma}

\begin{proposition}[Determinant of an alternating similitude]
Let \(V\) be a \(2m\)-dimensional \(E\)-vector space, let
\(B:V\times V\to E\) be a perfect alternating pairing, let
\(F\in\operatorname{Aut}_E(V)\), and let \(\mu\in E^\times\) satisfy
\[
                         B(Fx,Fy)=\mu B(x,y)
                         \qquad(x,y\in V).
\]
Then
\[
 \det(F)=\mu^m.
\]
\deps{The symplectic similitude determinant identity, normalized by the
connected symplectic group or by a symplectic basis and Pfaffians.}
\end{proposition}

\begin{proposition}[Determinant ambiguity for a symmetric similitude]
Let \(V\) be a \(b\)-dimensional \(E\)-vector space, let
\(B:V\times V\to E\) be a perfect symmetric pairing, let
\(F\in\operatorname{Aut}_E(V)\), and let \(\mu\in E^\times\) satisfy
\[
                         B(Fx,Fy)=\mu B(x,y)
                         \qquad(x,y\in V).
\]
Then
\[
 \det(F)^2=\mu^b.
\]
If \(\rho\in E^\times\) is chosen with \(\rho^2=\mu^b\), then the normalized
determinant \(\det(F)\rho^{-1}\) has square \(1\), hence belongs to
\(\{1,-1\}\).
\deps{Taking determinants in the matrix identity
\(F^{\mathsf t}BF=\mu B\).}
\end{proposition}

\section{A bounded graded duality system}

\begin{definition}[Graded perfect-similitude system]
Fix \(d\geq0\).  A graded perfect-similitude system of dimension parameter
\(d\) over \(E\) consists of:
\begin{enumerate}[label=\textup{(\roman*)}]
\item finite-dimensional spaces \(V_i\) and automorphisms \(F_i\),
      for \(0\leq i\leq2d\);
\item perfect pairings
\[
 B_i:V_i\times V_{2d-i}\longrightarrow E
\]
      satisfying
\[
 B_i(F_ix,F_{2d-i}y)=q^dB_i(x,y);
\]
\item a middle pairing \(B_d\) with swap sign \(\,(-1)^d\).
\end{enumerate}
Define
\[
 b_i=\dim_EV_i,\qquad D_i=\det(F_i),\qquad
 P_i(T)=\det(1-TF_i\mid V_i),
\]
\[
 \chi=\sum_{i=0}^{2d}(-1)^ib_i,\qquad
 \mathcal D=\prod_{i=0}^{2d}D_i^{(-1)^i},
\]
\[
 Z_F(T)=\prod_{i=0}^{2d}P_i(T)^{(-1)^{i+1}}.
\]
\end{definition}

\begin{lemma}[Complementary dimensions and determinants]
For every \(0\leq i\leq2d\),
\[
 b_i=b_{2d-i},\qquad D_iD_{2d-i}=q^{db_i}.
\]
\deps{Perfectness of \(B_i\) and the determinant formula for a similitude
pair of multiplier \(q^d\).}
\end{lemma}

\begin{lemma}[Complementary reciprocal polynomials]
For every \(i\),
\[
 P_{2d-i}(q^{-d}T^{-1})
 =(-1)^{b_i}T^{-b_i}D_i^{-1}P_i(T).
\]
Equivalently,
\[
 P_{2d-i}(T)=(-1)^{b_i}q^{db_i}D_i^{-1}T^{b_i}
 P_i(q^{-d}T^{-1}).
\]
\deps{The degreewise reciprocal-polynomial identity with multiplier
\(q^d\).}
\end{lemma}

\begin{lemma}[Parity of the half exponent]
The integer \(d\chi\) is even.  Hence
\[
 h=\frac{d\chi}{2}\in\mathbf Z.
\]
\deps{If \(d\) is even the assertion is immediate; if \(d\) is odd,
complementary degrees occur in pairs and the middle dimension \(b_d\) is
even because the middle pairing is alternating.}
\end{lemma}

\begin{lemma}[Square of the alternating Frobenius determinant]
The scalar
\(\mathcal D=\prod_iD_i^{(-1)^i}\) satisfies
\[
 \mathcal D^2=q^{d\chi}=q^{2h}.
\]
\deps{Complementary-degree reindexing of finite products and the identities
\(D_iD_{2d-i}=q^{db_i}\).}
\end{lemma}

\begin{definition}[Middle determinant sign]
Put
\[
 \delta=D_dq^{-db_d/2}.
\]
The exponent \(db_d/2\) is an integer: this is automatic when \(d\) is
even, and \(b_d\) is even when \(d\) is odd.
\end{definition}

\begin{proposition}[Middle-degree sign]
The scalar \(\delta\) satisfies \(\delta^2=1\).  More precisely:
\begin{enumerate}[label=\textup{(\roman*)}]
\item if \(d\) is odd, the middle pairing is alternating and \(\delta=1\);
\item if \(d\) is even, the middle pairing is symmetric and
      \(\delta\in\{1,-1\}\).
\end{enumerate}
\deps{The alternating-similitude determinant formula in odd \(d\), and the
symmetric-similitude determinant formula in even \(d\).}
\end{proposition}

\begin{proposition}[Exact determinant constant]
With \(h=d\chi/2\),
\[
 \mathcal D=\delta^{(-1)^d}q^h.
\]
Consequently
\[
 \varepsilon=(-1)^\chi q^{-h}\mathcal D
 =(-1)^\chi\delta^{(-1)^d}
\]
belongs to \(\{1,-1\}\).
\deps{Complementary-degree determinant products, the isolated middle
determinant, and the middle-degree sign proposition.}
\end{proposition}

\section{Raw and normalized functional equations}

\begin{definition}[Alternating zeta exponents]
For \(0\leq i\leq2d\), put
\[
 e_i=(-1)^{i+1}\in\{1,-1\}.
\]
Thus \(e_{2d-i}=e_i\),
\(Z_F(T)=\prod_iP_i(T)^{e_i}\),
\(\chi=-\sum_i e_i b_i\), and
\(\mathcal D=\prod_iD_i^{-e_i}\).
\end{definition}

\begin{lemma}[Complementary-degree reindexing of the zeta factors]
In \(E(T)\),
\[
 Z_F(q^{-d}T^{-1})
 =\prod_{i=0}^{2d}
 P_{2d-i}(q^{-d}T^{-1})^{e_i}.
\]
\deps{The involution \(i\mapsto2d-i\) of the finite index set and the
identity \(e_{2d-i}=e_i\).}
\end{lemma}

\begin{lemma}[Substitution of the complementary reciprocal factors]
In \(E(T)\),
\[
 Z_F(q^{-d}T^{-1})
 =\prod_{i=0}^{2d}
 \bigl((-1)^{b_i}T^{-b_i}D_i^{-1}P_i(T)\bigr)^{e_i}.
\]
\deps{Complementary-degree reindexing and the reciprocal-polynomial
identity for each pair of complementary degrees.}
\end{lemma}

\begin{lemma}[Separation and reassembly of the characteristic factors]
In \(E(T)\),
\begin{align*}
 &\prod_{i=0}^{2d}
 \bigl((-1)^{b_i}T^{-b_i}D_i^{-1}P_i(T)\bigr)^{e_i}\\
 &\quad=
 \left(\prod_{i=0}^{2d}((-1)^{b_i})^{e_i}\right)
 \left(\prod_{i=0}^{2d}(T^{-b_i})^{e_i}\right)
 \left(\prod_{i=0}^{2d}(D_i^{-1})^{e_i}\right)Z_F(T).
\end{align*}
\deps{Commutativity of \(E(T)^\times\), the exponent laws for
\(e_i\in\{1,-1\}\), and the defining product for \(Z_F(T)\).}
\end{lemma}

\begin{lemma}[Total Laurent exponent]
In \(E(T)^\times\),
\[
 \prod_{i=0}^{2d}(T^{-b_i})^{e_i}=T^\chi.
\]
\deps{The exponent identity
\(-\sum_i e_i b_i=\sum_i(-1)^ib_i=\chi\).}
\end{lemma}

\begin{lemma}[Total sign]
In \(E^\times\),
\[
 \prod_{i=0}^{2d}((-1)^{b_i})^{e_i}=(-1)^\chi.
\]
\deps{Every sign is its own inverse and
\(\sum_i b_i\equiv\sum_i(-1)^ib_i=\chi\pmod2\).}
\end{lemma}

\begin{lemma}[Total determinant scalar]
In \(E^\times\),
\[
 \prod_{i=0}^{2d}(D_i^{-1})^{e_i}=\mathcal D.
\]
\deps{The equality \(-e_i=(-1)^i\) and the defining finite product for
\(\mathcal D\).}
\end{lemma}

\begin{proposition}[Raw graded functional equation]
For every graded perfect-similitude system,
\[
 \boxed{
 Z_F(q^{-d}T^{-1})
 =(-1)^\chi\mathcal D\,T^\chi Z_F(T)}
 \quad\text{in }E(T).
\]
\deps{Complementary-factor substitution, separation and reassembly of the
characteristic factors, the total Laurent exponent, the total sign, and the
total determinant scalar.}
\end{proposition}

\begin{theorem}[Weil-normalized functional equation]
Let \(h=d\chi/2\) and
\(\varepsilon=(-1)^\chi q^{-h}\mathcal D\in\{1,-1\}\).  Then
\[
 \boxed{
 Z_F(q^{-d}T^{-1})
 =\varepsilon q^{d\chi/2}T^\chi Z_F(T).}
\]
\deps{The raw graded functional equation and the exact determinant
constant.}
\end{theorem}

\begin{corollary}[Central-variable normalization]
In a coefficient extension containing a chosen \(q^{1/2}\), define
\[
 \widetilde Z_F(U)=Z_F(q^{-d/2}U).
\]
Then
\[
 \boxed{\widetilde Z_F(U^{-1})
 =\varepsilon U^\chi\widetilde Z_F(U).}
\]
\deps{The change of variables \(T=q^{-d/2}U\) and the identity
\(q^{d\chi/2}(q^{-d/2}U)^\chi=U^\chi\).}
\end{corollary}

\begin{corollary}[The sign in the two symmetry cases]
If \(d\) is odd, then \(\chi\) is even and \(\varepsilon=1\).  If \(d\) is
even, then
\[
 \varepsilon=(-1)^\chi\delta,
\]
where \(\delta\) is the determinant sign of the symmetric middle pairing.
\deps{The parity of the half exponent and the middle-degree sign formula.}
\end{corollary}

\section{Realization-independent functional-equation constants}

Assume a compatible weighted realization family is concentrated in degrees
\(0,\ldots,2d\), and that every realization carries the graded
perfect-similitude pairings of the preceding sections.

\begin{definition}[Common determinants]
For the common integral factor
\[
 P_i(T)=1+\cdots+c_{i,b_i}T^{b_i},
\]
define
\[
 D_i^{\mathrm{com}}=(-1)^{b_i}c_{i,b_i}\in\mathbf Z\setminus\{0\}.
\]
Define
\[
 \chi^{\mathrm{com}}=\sum_i(-1)^ib_i,
 \qquad
 \mathcal D^{\mathrm{com}}
 =\prod_i(D_i^{\mathrm{com}})^{(-1)^i}\in\mathbf Q^\times,
\]
\end{definition}

\begin{lemma}[Leading coefficient recovers every determinant]
For every realization \(\lambda\), scalar extension carries
\(D_i^{\mathrm{com}}\) to
\(\det(F_{i,\lambda})\).  Hence the common Euler characteristic and
alternating determinant map to the corresponding quantities in every
realization.
\deps{The leading coefficient formula for
\(\det(1-TF_{i,\lambda})\) and individual-factor independence.}
\end{lemma}

\begin{proposition}[Parity and determinant square over the common field]
One has
\[
 d\chi^{\mathrm{com}}\in2\mathbf Z,
 \qquad
 (\mathcal D^{\mathrm{com}})^2=q^{d\chi^{\mathrm{com}}},
\]
and these identities map to the corresponding parity and determinant-square
identities in every realization.
\deps{Leading-coefficient recovery, one realization's parity lemma, the
square of its alternating determinant, and injectivity of
characteristic-zero scalar extension on \(\mathbf Q\).}
\end{proposition}

\begin{definition}[Common half exponent and sign]
Let \(h^{\mathrm{com}}\in\mathbf Z\) be the unique integer satisfying
\[
 2h^{\mathrm{com}}=d\chi^{\mathrm{com}}.
\]
Define
\[
 \varepsilon^{\mathrm{com}}
 =(-1)^{\chi^{\mathrm{com}}}q^{-h^{\mathrm{com}}}
  \mathcal D^{\mathrm{com}}\in\mathbf Q^\times.
\]
\end{definition}

\begin{proposition}[Independence of the functional-equation sign]
One has \(\varepsilon^{\mathrm{com}}\in\{1,-1\}\), and every coefficient
realization has this same sign.
\deps{The common determinant-square identity, the definition of the half
exponent, leading-coefficient recovery, and injectivity of
characteristic-zero scalar extension on \(\mathbf Q\).}
\end{proposition}

\begin{theorem}[Functional equation over the common field]
Let
\[
 Z(T)=\prod_{i=0}^{2d}P_i(T)^{(-1)^{i+1}}
 \in\mathbf Q(T).
\]
Then
\[
 \boxed{
 Z(q^{-d}T^{-1})
 =\varepsilon^{\mathrm{com}}
   q^{h^{\mathrm{com}}}
   T^{\chi^{\mathrm{com}}}Z(T).}
\]
The exponent, determinant constant, and sign are independent of the
auxiliary realization.
\deps{The realization-wise normalized functional equation, common integral
factors, leading-coefficient recovery, and injectivity of scalar extension
of rational functions.}
\end{theorem}

\section{Final arithmetic theorem for a smooth proper scheme}

\begin{theorem}[Arithmetic conclusion from trace, weights, and duality]
Let \(X_0\) be a smooth proper scheme of pure dimension \(d\) over
\(\Fq\).  For every prime \(\ell\nmid q\), write \(F_\ell\) for geometric
Frobenius acting degreewise on
\(H^i(X_0\times_{\Fq}\Fbar,\Ql)\), and suppose the following statements
hold:
\begin{enumerate}[label=\textup{(\roman*)}]
\item the cohomology spaces are finite-dimensional, vanish outside
      \(0\leq i\leq2d\), and Frobenius acts invertibly;
\item the Grothendieck--Lefschetz trace formula holds for every positive
      iterate;
\item after scalar extension to \(\Qlb\), every Frobenius eigenvalue in
      degree \(i\) is universally pure of weight \(i\);
\item Poincare duality supplies perfect pairings between degrees \(i\) and
      \(2d-i\), with multiplier \(q^d\), and middle swap sign
      \(\,(-1)^d\).
\end{enumerate}
For each such prime, let
\(\iota_\ell:\mathbf Z[T]\to\Ql[T]\) denote coefficientwise extension.
Then there are unique polynomials
\[
 P_i(T)\in\mathbf Z[T],\qquad P_i(0)=1,
 \qquad0\leq i\leq2d,
\]
such that, for every \(\ell\nmid q\),
\[
 \iota_\ell(P_i)(T)
 =\det(1-TF_\ell\mid H^i(X_0\times_{\Fq}\Fbar,\Ql)).
\]
They satisfy all of the following conclusions.
\begin{enumerate}[label=\textup{(\alph*)}]
\item \textup{Euler product and rationality:}
\[
 Z(X_0,T)=\prod_{x\in|X_0|}(1-T^{\deg x})^{-1}
 =\prod_{i=0}^{2d}P_i(T)^{(-1)^{i+1}}
 \in\mathbf Q(T).
\]
\item \textup{Integral factors and realization-independent degrees:}
\[
 b_i:=\deg P_i
 =\dim_{\Ql}H^i(X_0\times_{\Fq}\Fbar,\Ql)
\]
is independent of \(\ell\).
\item \textup{Riemann-hypothesis modulus:} every reciprocal root
      \(\alpha\) of \(P_i\) is an algebraic integer and
\[
 |\tau(\alpha)|=q^{i/2}
\]
      for every embedding
      \(\tau:\mathbf Q(\alpha)\hookrightarrow\mathbf C\).  Equivalently,
      every ordinary root \(\beta=\alpha^{-1}\) of \(P_i\) satisfies
      \( |\tau(\beta)|=q^{-i/2}\).
\item \textup{Poincare reciprocity:} if
      \(D_i=(-1)^{b_i}[T^{b_i}]P_i(T)\), then
\[
 P_{2d-i}(q^{-d}T^{-1})
 =(-1)^{b_i}T^{-b_i}D_i^{-1}P_i(T).
\]
      This is an equality in \(\mathbf Q[T,T^{-1}]\).
\item \textup{Functional equation:} with
\[
 \chi=\sum_i(-1)^ib_i,\qquad
 \varepsilon=(-1)^\chi q^{-d\chi/2}
 \prod_iD_i^{(-1)^i}\in\{1,-1\},
\]
      one has
\[
 \boxed{
 Z(X_0,q^{-d}T^{-1})
 =\varepsilon q^{d\chi/2}T^\chi Z(X_0,T).}
\]
\end{enumerate}
\deps{The closed-point Euler identity, cohomological determinant
factorization, Hankel rationality descent, Integral Fatou, no cancellation
between different weights, individual-factor recovery, perfect-similitude
reciprocity, and the Weil-normalized graded functional equation.}
\end{theorem}

\end{document}
